\section{The two zero motions and the corrected Hurwitz sign}
\label{sec:motion-local}

We compute the tangent vectors of the heat family and the Hurwitz translation in the same arithmetic coordinate. This identifies precisely which motion is measured by the phase statistic. The Hurwitz zero-velocity equation was already written by Frisk and de Gosson \cite[Section~2, equation~(3)]{frisk-degosson}; Garunk\v stis and Steuding studied its trajectories and their endpoint classifications \cite[Section~2]{garunkstis-steuding}. We combine the exact velocity with the midpoint phase convention below.
Write
\[
 \xi(s)=\tfrac12s(s-1)\pi^{-s/2}\Gamma(s/2)\zeta(s),
 \qquad H_0(x)=\tfrac18\xi(\tfrac12+ix/2).
\]
The heat family satisfies $\partial_tH_t=-\partial_x^2H_t$.
The second family is $F(q,s)=\zeta(s,1+q)$, where the last function is the Hurwitz zeta function. At $q=0$ it equals $\zeta(s)$. The identity
\[
 \partial_qF(q,s)=-s\zeta(s+1,1+q)
\]
follows first by differentiating the absolutely convergent series for $\Re s>1$ and then by meromorphic continuation; see \cite[25.11.17]{dlmf}.

\begin{proposition}[Exact tangent vectors]\label{prop:two-tangents}
At a simple real zero $x$ of $H_0$, the heat zero and its arithmetic coordinate satisfy
\[
 x'(0)=\frac{H_0''(x)}{H_0'(x)}\in\mathbb R,
 \qquad
 \left.\frac{d}{dt}\left(\tfrac12+ix(t)/2\right)\right|_{t=0}
       =\frac{iH_0''(x)}{2H_0'(x)}.
\]
The heat branch stays real for sufficiently small real $t$ of either sign, so its arithmetic coordinate has real part $1/2$ on that interval.
At a simple zero $\rho$ of $\zeta$, the Hurwitz branch has derivative
\[
 v_\rho:=s'(0)=\frac{\rho\zeta(\rho+1)}{\zeta'(\rho)}.
\]
\end{proposition}
\begin{proof}
The holomorphic implicit function theorem applies because the respective derivatives with respect to the zero coordinate are nonzero. Differentiating the two zero equations gives the displayed formulas. The function $H_t$ is real on the real axis for real $t$. Conjugating its unique local zero section produces another section with the same initial value; uniqueness identifies the two sections. Thus $x(t)$ is real throughout a sufficiently small real parameter interval. The affine map $x\mapsto1/2+ix/2$ gives the asserted arithmetic coordinate and its derivative.
\end{proof}

Let $\theta(u)$ denote the continuous Riemann--Siegel theta function with $\theta(0)=0$, and let
\[
 Z(u)=e^{i\theta(u)}\zeta(1/2+iu).
\]
At a positive zero ordinate, use the midpoint convention in
\[
 N(u)=\theta(u)/\pi+1+S(u).
\]
Under the zero-location input, enumerate the positive ordinates as $\gamma_n$, including every multiplicity. At a simple ordinate define
\[
 \phi_n=\arg\zeta(3/2+i\gamma_n),\qquad
 W_n=S(\gamma_n)-\phi_n/\pi,\qquad
 \delta_n=\arctan\frac1{2\gamma_n}.
\]
Here $\phi_n$ uses the logarithm specified by the absolutely convergent Euler product.

\begin{theorem}[The corrected Hurwitz sign]\label{thm:hurwitz-sign}
For every simple positive-ordinate zero $\rho_n=1/2+i\gamma_n$,
\[
 \Re v_{\rho_n}
 =\frac{|\zeta(\rho_n+1)|}{|Z'(\gamma_n)|}
       \left(\gamma_n\sin\pi W_n+\tfrac12\cos\pi W_n\right).
\]
In particular the strict rightward sign is exactly
$\sin(\pi W_n+\delta_n)>0$, and tangency occurs exactly when
$W_n+\delta_n/\pi\in\mathbb Z$.
At the indexed zero $n=1169$, the sine without $\delta_n$ is negative and the corrected sine is positive.
\end{theorem}
\begin{proof}
At a zero, differentiating $Z$ yields
$Z'(\gamma_n)=ie^{i\theta(\gamma_n)}\zeta'(\rho_n)$.
Since $Z(0)<0$, the sign at the crossing of the $n$th simple zero is
$\operatorname{sgn}Z'(\gamma_n)=(-1)^{n+1}$. This parity remains valid when preceding zeros have higher multiplicity: an even multiplicity changes neither the sign nor the parity of the count, while an odd multiplicity changes both.
The midpoint count reads
$\theta(\gamma_n)/\pi=n-3/2-S(\gamma_n)$. Hence
\[
 \frac1{\zeta'(\rho_n)}
 =\frac{ie^{i\theta(\gamma_n)}}{Z'(\gamma_n)}
 =\frac{e^{-i\pi S(\gamma_n)}}{|Z'(\gamma_n)|}.
\]
Insert $\zeta(\rho_n+1)=|\zeta(\rho_n+1)|e^{i\phi_n}$ into
Proposition~\ref{prop:two-tangents} and take real parts. Finally,
\[
 \gamma_n\sin\pi W_n+\tfrac12\cos\pi W_n
 =\sqrt{\gamma_n^2+1/4}\sin(\pi W_n+\delta_n),
\]
and the prefactor is strictly positive.

The interval computation described in the workbench uses the certified indexed zeta-zero routine and complex-ball power series for $\zeta'$. It gives
\begin{align*}
 \gamma_{1169}&=1613.0788554510274235\ldots,\\
 W_{1169}&=-0.0000606299724187599509\ldots,\\
 \delta_{1169}/\pi&=0.0000986653178528222938\ldots,\\
 \Re v_{\rho_{1169}}&=0.025925651967163814\ldots.
\end{align*}
The saved enclosures exclude zero for the last quantity and for each of the two sine values, with the signs stated in the theorem. The argument of the exterior value can be evaluated by its principal branch here: the Euler-product logarithm obeys
$|\Im\log\zeta(3/2+iu)|\le\log\zeta(3/2)<\pi$, so the two branches agree.
\end{proof}

\subsection{The exact periodic information required by the sign}
For every odd integer $k$, the midpoint convention gives
\[
 e^{i\pi kW_n}
 =e^{-3\pi ik/2}(-1)^n
       e^{-ik\theta(\gamma_n)}e^{-ik\phi_n}.
\]
Put $Y_n=W_n+\delta_n/\pi$. Its strict direction is
$g(Y_n)$, where $g(y)=\operatorname{sgn}(\sin\pi y)$ and
$g(y)=0$ at the integers. We give an exact finite estimate for the
mean of this sign, and a model that preserves Gaussian moments and
all even Fourier modes while keeping the sign positive.

\begin{proposition}[Periodic sign and Gaussian moments]
\label{prop:periodic-gaussian-gate}
There are real random variables $X_\sigma$, $\sigma\to\infty$,
whose normalized moments converge to every standard Gaussian moment,
and for which, for every nonzero integer $\ell$ and integer $k\ge0$,
\[
 \mathbb E\left[e^{2\pi i\ell X_\sigma}
                     (X_\sigma/\sigma)^k\right]\longrightarrow0,
 \qquad g(X_\sigma)=1\quad\hbox{almost surely}.
\]
For arbitrary real $y_1,\ldots,y_M$, $M\ge1$, put
$\langle f\rangle=M^{-1}\sum_{j=1}^M f(y_j)$ and
$b_\epsilon=M^{-1}\#\{j:\operatorname{dist}(y_j,\mathbb Z)<\epsilon\}$.
For $1/2\le r<1$, $0<\epsilon\le1/2$, and integer $J\ge0$,
\begin{align*}
 |\langle g\rangle|
 \le{}&\frac4\pi\sum_{l=0}^J
      \frac{r^{2l+1}}{2l+1}
       |\langle\sin((2l+1)\pi y)\rangle|+b_\epsilon\\
 &+\frac{3(1-r)}{\pi\sin(\pi\epsilon)}
 +\frac{4r^{2J+3}}{\pi(2J+3)(1-r^2)}.
\end{align*}
\end{proposition}
\begin{proof}
Let $G$ be standard normal and let $U$ be independent and uniform
on $(0,1)$. Set
\[
 X_\sigma=2\lfloor\sigma G/2\rfloor+U.
\]
Then $|X_\sigma-\sigma G|\le2$, so
$X_\sigma/\sigma\to G$ in every fixed $L^p$, $p>0$.
Also $g(X_\sigma)=\operatorname{sgn}(\sin\pi U)=1$.
Writing $A_\sigma=2\lfloor\sigma G/2\rfloor/\sigma$, the binomial
expansion of $(A_\sigma+U/\sigma)^k$ has leading term
$A_\sigma^k$. Its product with $e^{2\pi i\ell U}$ has expectation
zero, by independence and $\mathbb E e^{2\pi i\ell U}=0$.
Every remaining term has a factor $\sigma^{-1}$ or smaller,
and the moments of $A_\sigma$ are bounded. This proves all the
weighted limits. The exact map on periodic phases is
$\mathbb R/2\mathbb Z\to\mathbb R/\mathbb Z$, $y\mapsto y$:
the uniform measure on the half-interval $(0,1)$ pushes forward to
the uniform measure on the target circle. Its lift still has sign
one. Replacing $U$ by $U+1$ gives sign minus one with the same
limiting moments and even Fourier modes.

For the finite estimate, the absolutely convergent Abel sum is
\[
 g_r(y)=\frac4\pi\sum_{l\ge0}
       \frac{r^{2l+1}\sin((2l+1)\pi y)}{2l+1}
 =\frac2\pi\arctan\frac{2r\sin\pi y}{1-r^2}.
\]
The identity follows from the power series for
$\tfrac12\log((1+z)/(1-z))$ at $z=re^{i\pi y}$.
The relevant quotient has positive real part, which fixes the
arctangent branch. Hence $g_r$ has the sign of $g$, with
$|g_r|\le1$, and, away from the $\epsilon$-neighborhood of the
integers,
\[
 |g(y)-g_r(y)|
 =\frac2\pi\arctan\frac{1-r^2}{2r|\sin\pi y|}
 \le\frac{3(1-r)}{\pi\sin(\pi\epsilon)}.
\]
On the remaining set the difference is at most one.
Finally, the omitted Fourier tail has absolute value at most
$4r^{2J+3}/[\pi(2J+3)(1-r^2)]$. Averaging, applying the triangle
inequality, and combining these three estimates proves the claim.
\end{proof}

For a finite collection of simple zeta zeros, apply the estimate with
$y_j=Y_{n_j}$. Its left side is exactly the absolute difference
between the fractions of strict rightward and strict leftward
branches. The displayed odd-frequency means and boundary counts
are therefore explicit unresolved quantities for half-balance.
The correction from $W_n$ to $Y_n$ changes a sign only inside a
neighborhood of an integer of radius $\delta_n/\pi$; its counting
cost is exactly bounded by the number in that neighborhood.
